%============================================================
% boxes.tex — every coloured box used in the book.
% All boxes are ENVIRONMENTS, so lstlisting/tikz work inside.
% Boxes with [optional] titles keep a sensible default.
%============================================================

% Common skeleton for big titled boxes (theory / pattern)
\tcbset{
    dpbig/.style={
        enhanced, breakable,
        coltext=textonlight,
        boxrule=2pt, arc=4mm,
        left=8pt, right=8pt, top=8pt, bottom=6pt,
        drop shadow={shadow xshift=2pt, shadow yshift=-2pt, opacity=0.25},
        fonttitle=\bfseries\Large, coltitle=textondark,
        attach boxed title to top left={yshift=-3mm, xshift=5mm},
        before skip=10pt, after skip=10pt
    },
    dpsmall/.style={
        enhanced, breakable,
        coltext=textonlight,
        boxrule=1.5pt, arc=3mm,
        left=6pt, right=6pt, top=4pt, bottom=4pt,
        fonttitle=\bfseries, coltitle=textondark,
        before skip=8pt, after skip=8pt
    }
}

% --- Chapter-level theory (blue) ---
\newtcolorbox{chaptertheorybox}[1]{dpbig,
    colback=theorybg, colframe=theoryborder, colbacktitle=theoryborder,
    title={#1},
    boxed title style={enhanced, colframe=theoryborder, boxrule=0pt, arc=2mm}}

% --- Pattern-level theory (gold) ---
\newtcolorbox{patternbox}[1]{dpbig,
    colback=patternbg, colframe=patternborder, colbacktitle=patternborder,
    title={#1},
    boxed title style={enhanced, colframe=patternborder, boxrule=0pt, arc=2mm}}

% --- Solution flow of every solved problem ---
%   modelbox        what is computed over which elements, ranges or pairs, and the sizes that fix the budget
%   bruteforcebox   Approach 1: the naive solution and why it is too slow
%   algorithmbox    Approach 2: the efficient algorithm (code + complexity)
%   alternativebox  Approach 3: a second correct algorithm (other traversal, other structure)
%   tricksbox       Approach 4: implementation refinement (memory, constant factor, iterative form)
%   tracebox        step-by-step run of the algorithm on a small array
%   invariantbox    the invariant / lemma that proves the algorithm correct
\newtcolorbox{modelbox}[1][Problem Model]{dpsmall,
    colback=definitionbg, colframe=definitionborder, colbacktitle=definitionborder, title={#1}}
\newtcolorbox{bruteforcebox}[1][Approach 1: Brute Force]{dpsmall,
    colback=recursionbg, colframe=recursionred, colbacktitle=recursionred, title={#1}}
\newtcolorbox{algorithmbox}[1][Approach 2: Efficient Algorithm]{dpsmall,
    colback=tabulationbg, colframe=tabulationblue, colbacktitle=tabulationblue, title={#1}}
\newtcolorbox{alternativebox}[1][Approach 3: Alternative Algorithm]{dpsmall,
    colback=memobg, colframe=memoorange, colbacktitle=memoorange, title={#1}}
\newtcolorbox{tricksbox}[1][Approach 4: Implementation Refinements]{dpsmall,
    colback=optimizedbg, colframe=optimizedpurple, colbacktitle=optimizedpurple, title={#1}}
\newtcolorbox{tracebox}[1][Algorithm Trace]{dpsmall,
    colback=white, colframe=secondaryblue, colbacktitle=secondaryblue, title={#1}}
\newtcolorbox{invariantbox}[1][Invariant]{dpsmall,
    colback=white, colframe=deepblue, colbacktitle=deepblue, title={#1}}

% --- Planned chapter (stub shown until the chapter is written) ---
\newtcolorbox{plannedbox}[1][Planned Contents]{dpsmall,
    colback=lightgray!60, colframe=midgray, colbacktitle=darkgray, title={#1}}

% --- The four DP approaches (kept for the few DP-flavoured sections) ---
\newtcolorbox{recursionbox}[1][Approach 1: Pure Recursion (Brute Force)]{dpsmall,
    colback=recursionbg, colframe=recursionred, colbacktitle=recursionred, title={#1}}
\newtcolorbox{memobox}[1][Approach 2: Memoization (Top-Down DP)]{dpsmall,
    colback=memobg, colframe=memoorange, colbacktitle=memoorange, title={#1}}
\newtcolorbox{tabulationbox}[1][Approach 3: Tabulation (Bottom-Up DP)]{dpsmall,
    colback=tabulationbg, colframe=tabulationblue, colbacktitle=tabulationblue, title={#1}}
\newtcolorbox{optimizedbox}[1][Approach 4: Space-Optimized DP]{dpsmall,
    colback=optimizedbg, colframe=optimizedpurple, colbacktitle=optimizedpurple, title={#1}}

% --- Highlights ---
\newtcolorbox{insightbox}[1][Key Insight]{dpsmall,
    colback=insightbg, colframe=insightborder, colbacktitle=insightborder, title={#1}}
\newtcolorbox{warningbox}[1][Common Mistake / Warning]{dpsmall,
    colback=warningbg, colframe=warningborder, colbacktitle=warningborder, title={#1}}
\newtcolorbox{definitionbox}[1][Definition]{dpsmall,
    colback=definitionbg, colframe=definitionborder, colbacktitle=definitionborder, title={#1}}
\newtcolorbox{theorembox}[1][Theorem]{dpsmall,
    colback=white, colframe=deepblue, colbacktitle=deepblue, title={#1}}
\newtcolorbox{examplebox}[1][Worked Example]{dpsmall,
    colback=examplebg, colframe=exampleborder, colbacktitle=exampleborder, title={#1}}
\newtcolorbox{complexitybox}[1][Complexity Analysis]{dpsmall,
    colback=lightgray, colframe=darkgray, colbacktitle=darkgray, title={#1}}
\newtcolorbox{dptablebox}[1][DP Table Visualization]{dpsmall,
    colback=white, colframe=secondaryblue, colbacktitle=secondaryblue, title={#1}}
\newtcolorbox{comparisonbox}[1][Approach Comparison]{dpsmall,
    colback=white, colframe=softpurple, colbacktitle=softpurple, title={#1}}
\newtcolorbox{checklistbox}[1][Checklist]{dpsmall,
    colback=white, colframe=primaryblue, colbacktitle=primaryblue, title={#1}}
\newtcolorbox{testbox}[1][Test Suite]{dpsmall,
    colback=white, colframe=successgreen, colbacktitle=successgreen, title={#1}}

% --- Untitled boxes ---
\newtcolorbox{recurrencebox}{enhanced,
    colback=white, colframe=deepblue, coltext=textonlight,
    boxrule=2pt, arc=3mm, left=10pt, right=10pt, top=6pt, bottom=6pt,
    before skip=8pt, after skip=8pt}
\newtcolorbox{diagrambox}[1][]{enhanced, breakable,
    colback=white, colframe=darkgray!40, coltext=textonlight,
    boxrule=0.5pt, arc=2mm, left=4pt, right=4pt, top=4pt, bottom=4pt,
    halign=center, before skip=8pt, after skip=8pt, #1}

%------------------------------------------------------------
% Problem header for SOLVED problems
%   \problemcard{id}{title}{source}{url}{difficulty}{technique}
%------------------------------------------------------------
\newcommand{\problemcard}[6]{%
\phantomsection\label{pb:#4}% anchor for the problem index (appendix)
\begin{tcolorbox}[enhanced, breakable,
    colback=lightgray, colframe=primaryblue, boxrule=2pt, arc=4mm,
    drop shadow={shadow xshift=2pt, shadow yshift=-2pt, opacity=0.3},
    title={Problem #1: #2}, fonttitle=\bfseries\large,
    coltitle=textondark, colbacktitle=primaryblue,
    attach boxed title to top left={yshift=-3mm, xshift=5mm},
    boxed title style={enhanced, colframe=primaryblue, boxrule=0pt, arc=2mm},
    top=8pt]
\renewcommand{\arraystretch}{1.25}%
\begin{tabularx}{\linewidth}{@{}>{\bfseries\color{deepblue}}l X@{}}
Source     & #3 \quad \solveit{#4} \\
Difficulty & #5 \\
Technique  & #6 \\
\end{tabularx}
\end{tcolorbox}}

%------------------------------------------------------------
% Compact entry for PRACTICE problem sets
%   \problembox{id}{title}{technique / hint}{url}
%------------------------------------------------------------
\newcommand{\problembox}[4]{%
\phantomsection\label{pb:#4}% anchor for the problem index (appendix)
\begin{tcolorbox}[enhanced, breakable,
    colback=white, colframe=primaryblue!70, boxrule=1pt, arc=2mm,
    left=6pt, right=6pt, top=3pt, bottom=3pt, before skip=4pt, after skip=4pt]
\textbf{\textcolor{deepblue}{#1\quad #2}}\\[2pt]
{\small\textcolor{darkgray}{#3}}\\[-1pt]
{\scriptsize\textcolor{secondaryblue}{\url{#4}}}\hfill\solveit{#4}
\end{tcolorbox}}

%------------------------------------------------------------
% Solution box
%------------------------------------------------------------
\newtcolorbox{solutionbox}[1][Solution]{dpsmall,
    colback=solutionbg, colframe=solutiongreen, colbacktitle=solutiongreen, title={#1}}

%------------------------------------------------------------
% Variants of a solved problem (end of every solved problem)
%   \begin{variantsbox}
%     \variant{<short name>}{<where>}{<what changes, 1 line>}{<subproblem, <= 2 lines>}{<recurrence + answer, <= 2 lines>}
%   \end{variantsbox}
%   <where> is one of
%     \exmedium\ \lc{Medium}\ \varlink{<judge title>}{<source>}{<url>}   problem NOT elsewhere in the book
%     \varin{<url>}                                                     problem solved elsewhere in the book
%     \varnone                                                          no online judge
%------------------------------------------------------------
\newcounter{variantnum}
\newtcolorbox{variantsbox}[1][Variants --- One Idea, Many Problems]{dpsmall,
    colback=softpurple!4, colframe=softpurple, colbacktitle=softpurple, title={#1},
    before upper={\setcounter{variantnum}{0}\small}}
\newcommand{\variant}[5]{%
  \stepcounter{variantnum}%
  \par\medskip\noindent{\bfseries\textcolor{softpurple!75!black}{V\arabic{variantnum}}\enspace #1}%
  {\unskip\nobreak\hfil\penalty50\hskip1em\hbox{}\nobreak\hfill{\footnotesize #2}\parfillskip=0pt\par}\vspace{2pt}%
  \noindent\hangindent=1.6em\hangafter=0\textcolor{darkgray}{\textit{Change:}}\enspace #3\par
  \noindent\hangindent=1.6em\hangafter=0\textcolor{darkgray}{\textit{Subproblem:}}\enspace #4\par
  \noindent\hangindent=1.6em\hangafter=0\textcolor{darkgray}{\textit{Recurrence:}}\enspace #5\par}
\newcommand{\varlink}[3]{\phantomsection\label{pb:#3}\href{#3}{\textcolor{deepblue}{#2}}~\href{#3}{\exbadge{primaryblue}{Solve $\nearrow$}}}
\newcommand{\varin}[1]{\hyperref[pb:#1]{\textcolor{deepblue}{in this book, p.~\pageref*{pb:#1}}}}
\newcommand{\varnone}{\textcolor{darkgray}{(no online judge)}}
